Each correction has components, and the ablation removes or changes them one at a time on the same replications, so that the contribution of each is a difference of paired rates. \Cref{tab:coverage-ablation}, \cref{tab:shift-ablation} and \cref{tab:situation-ablation} give the variants of the coverage test, the resilience test and the situation light.

\subsection{Components of the coverage test}

\begin{table}[htbp]
\centering
\small
\caption{The components of the coverage test removed or changed one at a time on the same covered counts: the size and the power at two shapes of drift for equal windows (the first three columns) and for a test window four times the calibration window (the last three); simulated rates with their standard errors.}
\label{tab:coverage-ablation}
\fitwidth{\begin{tabular}{p{5.2cm}cccccc}
\toprule
Variant & Size & Power, $b=1.05$ & Power, $b=1.10$ & Size  & Power, $b=1.05$  & Power, $b=1.10$  \\
\midrule
binomial p-value (published practice) & $0.157 \pm 0.006$ & $0.204 \pm 0.006$ & $0.372 \pm 0.008$ & $0.361 \pm 0.008$ & $0.412 \pm 0.008$ & $0.532 \pm 0.008$ \\
Beta-Binomial p-value, minimum-likelihood rule (proposed) & $0.051 \pm 0.003$ & $0.076 \pm 0.004$ & $0.190 \pm 0.006$ & $0.048 \pm 0.003$ & $0.079 \pm 0.004$ & $0.162 \pm 0.006$ \\
Beta-Binomial, equal-tail p-value & $0.052 \pm 0.004$ & $0.069 \pm 0.004$ & $0.179 \pm 0.006$ & $0.046 \pm 0.003$ & $0.070 \pm 0.004$ & $0.146 \pm 0.006$ \\
normal test with the variance ratio of Proposition 2 & $0.051 \pm 0.003$ & $0.076 \pm 0.004$ & $0.190 \pm 0.006$ & $0.047 \pm 0.003$ & $0.079 \pm 0.004$ & $0.162 \pm 0.006$ \\
Beta-Binomial with the calibration size halved & $0.018 \pm 0.002$ & $0.033 \pm 0.003$ & $0.102 \pm 0.005$ & $0.007 \pm 0.001$ & $0.020 \pm 0.002$ & $0.055 \pm 0.004$ \\
Beta-Binomial with the calibration size doubled & $0.089 \pm 0.005$ & $0.122 \pm 0.005$ & $0.264 \pm 0.007$ & $0.124 \pm 0.005$ & $0.164 \pm 0.006$ & $0.286 \pm 0.007$ \\
binomial with the exact size-calibrated cut & $0.052 \pm 0.004$ & $0.069 \pm 0.004$ & $0.179 \pm 0.006$ & $0.048 \pm 0.003$ & $0.072 \pm 0.004$ & $0.151 \pm 0.006$ \\
\bottomrule
\end{tabular}}
\end{table}


The first component is the law. Replacing the binomial law by the Beta-Binomial law lowers the size from \CovSinBinB{} to \CovSinBBB{} at equal windows and from \CovSinBinC{} to \CovSinBBC{} at four times the calibration window (\cref{tab:coverage-ablation}), and it lowers the power at the drift $b=1.10$ from \CovPowBinBTen{} to \CovPowBBBTen{} and from \CovPowBinCTen{} to \CovPowBBCTen{}. The second component is the rule that turns the law into a p-value. The equal-tail rule, which doubles the smaller tail, has the size \CovSinEqB{} and \CovSinEqC{} and the power \CovPowEqBTen{} and \CovPowEqCTen{} at $b=1.10$; the minimum-likelihood rule has a little more power (\CovPowBBBTen{} and \CovPowBBCTen{}), by a difference of about one to two standard errors, which does not separate the two rules. The third component is the exactness of the law. A normal test whose variance is the binomial variance times the factor of \cref{prop:variance} has the size \CovSinZinfB{} and \CovSinZinfC{} and the power \CovPowZinfBTen{} and \CovPowZinfCTen{}, indistinguishable from the exact law at these window sizes: the whole correction is carried by the variance factor, and the exact law matters for small windows, where the normal approximation of \cref{tab:exact-size} is off by as much as \CovNormalMaxErr{}.

The fourth component is the knowledge of the calibration size, which the law needs and the binomial test does not. A test that reads the law with half the true calibration size is conservative, with the size \CovSinHalfB{} and \CovSinHalfC{} and the power \CovPowHalfBTen{} and \CovPowHalfCTen{}; a test that reads it with double the size is oversized, \CovSinDoubleB{} and \CovSinDoubleC{}, though less than the binomial test at the same windows. An error of one rank in the quantile is nearly harmless: a test of rank $k$ applied to a calibration whose true rank is $k+1$ has the exact size \RankOffBBOne{} and \RankOffBBFour{} at the two windows, against \RankOffBinOne{} and \RankOffBinFour{} for the binomial test. The correction therefore needs the right calibration size, and the value of $n$ must be recorded with the conformal set, as the engine does; it does not need the rank to be right to the last unit. The fifth variant, the binomial test with its cut moved until the exact size is $0.05$, needs a p-value cut of \CovOracleCutB{} at equal windows and \CovOracleCutC{} at four times the calibration window instead of $0.05$; it has the size \CovSinOracleB{} and \CovSinOracleC{} and the power \CovPowOracleBTen{} and \CovPowOracleCTen{}, the power of the Beta-Binomial test within its error. A binomial test can be corrected by a cut that depends on $n$ and $m$, and the Beta-Binomial test gives the same behaviour without a cut that has to be computed for each pair $(n,m)$.

\subsection{Components of the resilience test}

\begin{table}[htbp]
\centering
\small
\caption{The rejection rate of each rule of the resilience test with 500 evaluation rows, 1,000 development rows and the worst quarter as the subsample: the size without shift and the power at three planted shifts, with the permutation reference read at four numbers of draws; rates with their standard errors.}
\label{tab:shift-ablation}
\fitwidth{\begin{tabular}{lcccc}
\toprule
Rule & Size (shift 0) & Power (shift 1) & Power (shift 2) & Power (shift 3) \\
\midrule
fixed drop above 0.05 & $0.175 \pm 0.019$ & $0.552 \pm 0.025$ & $0.925 \pm 0.013$ & $0.993 \pm 0.004$ \\
fixed drop above 0.10 & $0.048 \pm 0.011$ & $0.235 \pm 0.021$ & $0.703 \pm 0.023$ & $0.927 \pm 0.013$ \\
variance z rule & $0.065 \pm 0.012$ & $0.223 \pm 0.021$ & $0.723 \pm 0.022$ & $0.965 \pm 0.009$ \\
permutation reference, 19 draws & $0.000 \pm 0.000$ & $0.000 \pm 0.000$ & $0.000 \pm 0.000$ & $0.000 \pm 0.000$ \\
permutation reference, 39 draws & $0.037 \pm 0.009$ & $0.150 \pm 0.018$ & $0.463 \pm 0.025$ & $0.690 \pm 0.023$ \\
permutation reference, 99 draws & $0.055 \pm 0.011$ & $0.242 \pm 0.021$ & $0.618 \pm 0.024$ & $0.877 \pm 0.016$ \\
permutation reference, 199 draws & $0.062 \pm 0.012$ & $0.250 \pm 0.022$ & $0.680 \pm 0.023$ & $0.900 \pm 0.015$ \\
fixed drop above the 95th percentile of the null drops & $0.050 \pm 0.011$ & $0.242 \pm 0.021$ & $0.720 \pm 0.022$ & $0.927 \pm 0.013$ \\
\bottomrule
\end{tabular}}
\end{table}


\Cref{tab:shift-ablation} varies the rule at 500 evaluation rows. Standardising the drop with the variance of the structural components of the area, the z rule, brings the size from \ShRateFixedFiveCZ{} to \ShRateZruleCZ{} and keeps a power of \ShRateZruleCT{} at shift 2, close to the \ShRateFixedTenCT{} of the cut 0.10 and above the \ShRatePermCT{} of the reference. The reference does not use a closed-form variance; it measures the null law of the drop with the model, the selection of the worst subsample and the resemblance score all inside it, so it does not depend on the normal approximation of the area, which is the weaker the smaller the worst subsample. At 125 and 250 rows, where the worst quarter holds 32 and 63 rows, the reference has more power than the z rule (\cref{tab:shift-eval-size}), and at 500 and 1,000 rows less. The second component is the number of draws. With 39 draws the power at shift 2 is \ShRatePermThirtyNineCT{}, with 99 it is \ShRatePermNinetyNineCT{} and with 199 it is \ShRatePermCT{}, and the size rises from \ShRatePermThirtyNineCZ{} to \ShRatePermNinetyNineCZ{} and \ShRatePermCZ{} toward the level, as \cref{prop:permutation} predicts for a p-value grid that gets finer. The third component is the calibration of the cut by the null: the fixed cut at the 95th percentile of the null drops, the cut that needs a simulation the validation team does not have, has the size \ShOracleRateZ{} and the powers \ShOracleRateO{}, \ShOracleRateT{} and \ShOracleRateH{} at the shifts 1, 2 and 3, against \ShRatePermCO{}, \ShRatePermCT{} and \ShRatePermCH{} for the reference with 199 draws. The reference therefore reaches, within the error, the operating point of a cut that is calibrated on the null of the same cell, without the cell's null being simulated beforehand.

\subsection{Components of the situation-testing light}

\begin{table}[htbp]
\centering
\small
\caption{The rejection rate of nine rules of the situation-testing light: the size (gap 0) and the power at two planted gaps for 400 and for 1,000 applicants, with the standard error of every rate at most 0.025 over 1,000 replications at gap 0 and 400 at the others.}
\label{tab:situation-ablation}
\fitwidth{\begin{tabular}{lcccccc}
\toprule
Rule & $n=400$, gap 0 & gap 0.2 & gap 0.4 & $n=1{,}000$, gap 0 & gap 0.2 & gap 0.4 \\
\midrule
share cut 0.05 & 0.274 & 0.905 & 1.000 & 0.145 & 1.000 & 1.000 \\
share cut 0.075 & 0.170 & 0.833 & 1.000 & 0.050 & 0.980 & 1.000 \\
share cut 0.10 & 0.114 & 0.743 & 1.000 & 0.018 & 0.927 & 1.000 \\
share cut 0.15 & 0.052 & 0.595 & 0.988 & 0.000 & 0.740 & 1.000 \\
share cut 0.20 & 0.014 & 0.410 & 0.968 & 0.000 & 0.468 & 1.000 \\
share cut 0.25 & 0.005 & 0.268 & 0.932 & 0.000 & 0.190 & 1.000 \\
share cut 0.30 & 0.002 & 0.152 & 0.873 & 0.000 & 0.055 & 0.995 \\
binomial test of the count & 0.156 & 0.823 & 1.000 & 0.070 & 0.990 & 1.000 \\
one significant at the Bonferroni level & 0.079 & 0.453 & 0.938 & 0.053 & 0.760 & 1.000 \\
\bottomrule
\end{tabular}}
\end{table}


\Cref{tab:situation-ablation} varies the cut and the form of the rule. Moving the cut from 0.05 to 0.30 moves the size at 400 applicants from \SiRateCutFiveSmallZ{} to \SiRateCutThirtySmallZ{} and the power at the gap 0.2 from \SiRateCutFiveSmallB{} to \SiRateCutThirtySmallB{}; at 1,000 applicants the size is \SiRateCutFiveLargeZ{} at the cut 0.05 and zero from the cut 0.15, and the power at the gap 0.2 is \SiRateCutFiveLargeB{} and \SiRateCutThirtyLargeB{} at the two ends. The cut is the only parameter of the rule, so these are points on one operating curve, and the ablation of the cut is the curve. The binomial test of the count, which treats the complainants as independent, sits above the share cuts on that curve at 400 applicants: its size is \SiRateCountBinSmallZ{} and its power at the gap 0.2 is \SiRateCountBinSmallB{}, and the share cut 0.10 has the size \SiRateCutTenSmallZ{} and the power \SiRateCutTenSmallB{}, so that the binomial test buys its power with a size above the level, as the independence assumption predicts. The rule that fires when one complainant is significant at the Bonferroni level has the size \SiRateBonfSmallZ{} and the power \SiRateBonfSmallB{} at the same gap, and at 1,000 applicants the size \SiRateBonfLargeZ{} and the power \SiRateBonfLargeB{}. On this plant the share cuts dominate it: at 400 applicants the cut 0.15 has the smaller size (\SiRateCutFifteenSmallZ{}) and the larger power (\SiRateCutFifteenSmallB{}), and at 1,000 applicants the cut 0.10 has the size \SiRateCutTenLargeZ{} and the power \SiRateCutTenLargeB{}. The cut 0.20 is not on the efficient part of the curve at 1,000 applicants: its size is \SiRateCutTwentyLargeZ{}, like that of the cut 0.15, and its power at the gap 0.2 is \SiRateCutTwentyLargeB{} against \SiRateCutFifteenLargeB{}. The choice among these points is a choice of the error to avoid. The rule of \cref{eq:cut-rule} selects \SiChosenThrSmall{} on the threshold-bias plant at 400 applicants and \SiChosenThrLarge{} at 1,000, and \SiChosenIndSmall{} and \SiChosenIndLarge{} on the plant with independent attributes; the engine's cut of 0.20 is the choice of the second plant at 400 applicants, and it is conservative on the others.
